\clearpage
\section{Guía Rápida: Qué Usar y Cuándo}

Resumen en lenguaje llano de cada técnica del libro: qué resuelve, cuándo conviene
sobre las alternativas, y un ejemplo mínimo de entrada y salida. Sirve para el
momento de leer un enunciado y decidir por dónde atacarlo; el desarrollo completo
está en la página indicada. Los patrones de lectura rápida están en la sección~1 y
las utilidades de Python en la sección~2.

\newcommand{\guia}[2]{\medskip\noindent\textbf{#1} \textnormal{(p.~\pageref{#2})}\ }
\newcommand{\ej}{\par\smallskip\noindent\textit{Ejemplo:}\ }

% ============================================================
\subsection*{Estructuras de Datos}

\guia{Fenwick Tree (BIT)}{alg:3-1}
Un arreglo al que le vas a preguntar muchas veces "¿cuánto suman las posiciones de
aquí a allá?" mientras además vas cambiando valores. Recalcular la suma cada vez
sería demasiado lento; el BIT la responde en pasos logarítmicos. Es la estructura
más corta de escribir de todo el capítulo, así que es la primera opción cuando la
operación es una suma.
\ej con \texttt{[5, 2, 9, 1]}, sumar el rango \texttt{[2,4]} da \texttt{12}; si
cambias la posición 3 a \texttt{0}, la misma consulta pasa a dar \texttt{3}.

\guia{Sumas de Prefijos + Mapa}{alg:3-2}
Para encontrar el tramo contiguo más largo que está "equilibrado": misma cantidad
de dos cosas, o que suma exactamente un valor. El truco es convertir cada elemento
en $+1$ o $-1$ y buscar dos puntos donde el acumulado valga lo mismo: todo lo que
hay entre ellos está balanceado. Una sola pasada, sin revisar todos los tramos.
\ej en \texttt{[+,+,-,-,+,-]} el tramo más largo con igual cantidad de \texttt{+}
y \texttt{-} es de longitud \texttt{6}.

\guia{Fenwick con XOR (Nim)}{alg:3-3}
El mismo BIT pero combinando con XOR en vez de suma, que es lo que hace falta para
juegos por turnos tipo Nim: hay varias pilas, cada jugador retira de una sola, y
pierde quien no puede mover. El XOR de los tamaños decide el ganador, y el BIT
permite responder eso para muchos rangos distintos mientras las pilas cambian.
\ej pilas \texttt{[1, 2, 3]}: el XOR da \texttt{0}, así que pierde quien empieza.

\guia{Segment Tree}{alg:3-4}
Lo mismo que el BIT pero para operaciones que no se pueden "deshacer", como el
mínimo o el máximo: ahí no sirve restar. Es la estructura por defecto para
consultas de rango sobre un arreglo que cambia. Más largo de escribir que el BIT,
así que si la operación es una suma, usa el BIT.
\ej con \texttt{[3, 1, 4, 1, 5]}, el mínimo de \texttt{[1,3]} es \texttt{1}; tras
cambiar la posición 2 a \texttt{9}, pasa a ser \texttt{3}.

\guia{Dos Fenwick (rango + rango)}{alg:3-5}
Cuando además de consultar sumas de un rango necesitas SUMAR un valor a todo un
rango de golpe. Se resuelve con dos BIT en vez de un segment tree con propagación
perezosa: mucho menos código y varias veces más rápido en Python.
\ej sobre \texttt{[0,0,0,0]}, sumar \texttt{3} a \texttt{[1,3]} y luego consultar
\texttt{[2,4]} da \texttt{6}.

\guia{Sparse Table}{alg:3-6}
Si el arreglo NUNCA cambia, esta tabla responde el mínimo (o máximo, o gcd) de
cualquier rango en tiempo constante, más rápido que un segment tree. El precio es
que no admite modificaciones. Regla simple: arreglo fijo, sparse table; arreglo que
cambia, segment tree.
\ej con \texttt{[7, 2, 8, 2, 4]}, el mínimo de \texttt{[3,5]} es \texttt{2}, y la
respuesta sale mirando solo dos casillas ya calculadas.

\guia{Pila Monótona}{alg:3-7}
Responde, para cada elemento, hasta dónde se extiende hacia los lados sin encontrar
uno menor. Con eso se resuelve toda la familia de "el rectángulo más grande que
cabe", "hasta cuándo este precio sigue siendo el más bajo" o "cuánto dura esta
racha". Una sola pasada, aunque parezca que hay que probar todos los pares.
\ej alturas \texttt{[2, 1, 5, 6, 2, 3]}: el rectángulo más grande tiene área
\texttt{10} (las alturas \texttt{5} y \texttt{6} con ancho 2).

\guia{Sumas de Prefijos 2D}{alg:3-8}
La versión en cuadrícula de las sumas de prefijos: precalculas una tabla y después
cualquier pregunta del tipo "cuántas cosas hay dentro de este rectángulo" se
responde sumando y restando cuatro valores. Su gemelo, el arreglo de diferencias,
hace lo contrario: aplica miles de sumas sobre rectángulos y al final reconstruye
la cuadrícula de una vez.
\ej en una grilla con árboles marcados, contar los de un rectángulo cuesta lo mismo
sin importar si el rectángulo tiene 4 celdas o un millón.

\guia{Ordered Set con Fenwick}{alg:3-9}
Python no tiene el \texttt{ordered\_set} de C++, y esta es la forma de suplirlo:
mantener un conjunto donde puedes insertar, borrar, preguntar "¿cuántos elementos
son menores que $x$?" y "¿cuál es el $k$-ésimo más pequeño?". Requiere conocer de
antemano todos los valores que aparecerán, para comprimirlos.
\ej tras insertar \texttt{10, 40, 20}, el 2.º menor es \texttt{20} y hay
\texttt{2} elementos menores que \texttt{40}.

\guia{Balance de Paréntesis por Bloques}{alg:3-10}
Para secuencias de paréntesis balanceadas donde hay que razonar qué parejas se
pueden formar sin romper el resto. La idea clave es que la secuencia se parte en
bloques independientes, y solo se pueden emparejar elementos del mismo bloque.
\ej en \texttt{()()}, el primer \texttt{(} no puede emparejarse con el último
\texttt{)}: están en bloques distintos.

\guia{Pila con Restricción}{alg:3-11}
Validar estructuras anidadas (cajas dentro de cajas, etiquetas, expresiones) donde
además cada contenedor impone una condición sobre lo que lleva dentro. La pila
lleva el control de qué está abierto, y cada cierre verifica la condición.
\ej \texttt{-9 -2 2 -3 3 9} es válido (2 y 3 caben en el 9), pero \texttt{-9 -9 9
9} no lo es.

\guia{Reducción con Pila}{alg:3-12}
Simplificar un recorrido eliminando los rodeos: si pasaste por un sitio y volviste,
todo lo del medio sobra. Lo importante es que al quitar un rodeo puede aparecer
otro más grande, así que hay que seguir revisando.
\ej la ruta \texttt{a!b!c!b!a!d} se reduce a \texttt{a!d}: primero desaparece
\texttt{c}, y eso deja \texttt{b} junto a \texttt{b}.

\guia{Consultas Offline}{alg:3-13}
Cuando todas las preguntas se conocen de antemano, se pueden responder en el orden
que más convenga y no en el que llegan. Para "cuántos elementos mayores que $k$ hay
en este tramo", ordenar las consultas por $k$ hace que una de las dos condiciones se
cumpla sola. Luego solo queda contar posiciones en un rango; en
Python, con prefijos por bloques es más rápido que con un Fenwick.
\ej en \texttt{[5, 1, 2, 3, 4]}, entre las posiciones 2 y 4 hay \texttt{2}
elementos mayores que 1.

% ============================================================
\subsection*{Grafos}

\guia{Emparejamiento Bipartito (Kuhn)}{alg:4-1}
Dos grupos y una lista de parejas posibles entre ellos; se busca formar la mayor
cantidad de parejas sin repetir a nadie. Asignar trabajadores a tareas, estudiantes
a proyectos, turnos a personas. El algoritmo reacomoda parejas ya formadas cuando
eso permite meter una más.
\ej con 3 chicos, 3 chicas y las parejas posibles $1$--$1$, $1$--$2$, $2$--$1$,
la respuesta es \texttt{2}.

\guia{Kuhn sobre Segmentos}{alg:4-2}
Colocar piezas en un tablero con obstáculos sin que se ataquen. En lugar de pensar
celda por celda, se numeran los tramos libres de cada fila y de cada columna, y el
problema se vuelve el emparejamiento de arriba. Es el ejemplo típico de "modelar
para poder aplicar un algoritmo conocido".
\ej en un tablero $3\times3$ con un peón en el centro caben \texttt{4} torres sin
que se ataquen, en vez de las 3 del tablero vacío.

\guia{Teorema de König}{alg:4-3}
Cuántos elementos hay que quitar, como mínimo, para que no quede ningún conflicto.
En general es un problema durísimo, pero si los conflictos forman un grafo
bipartito la respuesta es exactamente el emparejamiento máximo: se calcula igual
que la entrada anterior y ya está.
\ej con caballos de ajedrez que se atacan, basta correr Kuhn; el número que sale
es la cantidad mínima a retirar.

\guia{Componentes Fuertemente Conexas}{alg:4-4}
En un grafo de calles de un solo sentido, agrupa los nodos entre los que se puede
ir y volver. Sirve para simplificar el grafo (cada grupo pasa a ser un solo nodo,
y lo que queda ya no tiene ciclos) y es el motor de 2-SAT.
\ej con $1\to2$, $2\to3$, $3\to1$ y $3\to4$, los planetas 1, 2 y 3 son un mismo
reino y el 4 es otro.

\guia{Circuito Euleriano (Hierholzer)}{alg:4-5}
Recorrer TODAS las calles exactamente una vez y volver al punto de partida. Ojo con
la confusión frecuente: esto es sobre aristas, no sobre nodos (pasar por todos los
nodos una vez es otro problema, mucho más difícil). Se comprueba antes si existe:
todos los nodos deben tener grado par.
\ej un triángulo tiene circuito (\texttt{1 2 3 1}); una "Y" de tres brazos no
tiene, porque sus cuatro nodos quedan con grado impar.

\guia{Puntos de Articulación (Tarjan)}{alg:4-6}
Qué nodos son críticos: si ese nodo falla, la red se parte en pedazos. Se usa para
medir robustez, encontrar cuellos de botella o identificar el servidor que no puede
caerse.
\ej en la cadena $1-2-3$, el nodo \texttt{2} es punto de articulación; en un
triángulo no hay ninguno.

\guia{Bron--Kerbosch (clique máximo)}{alg:4-7}
El grupo más grande en el que todos son compatibles con todos. Es exponencial, así
que solo sirve con pocos nodos (unas decenas), pero para ese tamaño es la
herramienta correcta y no hay atajo mejor.
\ej si A, B y C son compatibles entre sí y D solo con A, el grupo máximo es
\texttt{\{A, B, C\}}, de tamaño \texttt{3}.

\guia{Orden Topológico + Conteos}{alg:4-8}
Cuando entran cantidades astronómicas de elementos por un grafo sin ciclos y
simularlos uno por uno es imposible. En vez de mover elementos, se mueven CONTEOS:
se procesa cada nodo cuando ya sabe cuántos le llegaron en total, y reparte de una
sola vez.
\ej con \ten{18} bolas entrando por el nodo 1, no se simula ninguna bola: se
calcula de una cuántas pasan por cada nodo.

\guia{Union-Find (DSU)}{alg:4-9}
Va uniendo elementos en grupos y responde al instante si dos están en el mismo
grupo. Es la forma más rápida de contar cuántos grupos separados hay cuando las
conexiones van llegando una a una, y no necesita construir el grafo.
\ej con las uniones \texttt{A-B} y \texttt{C-D} sobre cuatro elementos quedan
\texttt{2} grupos.

\guia{Dijkstra}{alg:4-10}
La ruta más barata desde un punto a todos los demás, cuando ningún tramo tiene
costo negativo. Es el algoritmo de caminos mínimos por defecto: mapas, costos,
cualquier grafo donde avanzar cuesta. Si hay costos negativos, no sirve: ahí va
Bellman-Ford.
\ej con $1\to2$ costo 5 y $1\to3\to2$ costos 1 y 1, la distancia de 1 a 2 es
\texttt{2}, no 5.

\guia{BFS en grilla (y 0-1 BFS)}{alg:4-11}
El camino más corto cuando todos los pasos cuestan igual: laberintos, mínimo número
de movimientos, distancia en casillas. Recorre por capas, así que la primera vez
que llega a un sitio ya es por el camino más corto. La variante 0-1 permite que
algunos pasos sean gratis y otros cuesten uno, sin necesitar Dijkstra.
\ej en un laberinto, devuelve la longitud mínima y la secuencia de movimientos
\texttt{DRRU...} que hay que seguir.

\guia{Bipartito por coloreo}{alg:4-12}
¿Se puede partir el grupo en dos equipos de forma que ningún par conflictivo quede
junto? Se pinta cada nodo del color contrario al de su vecino; si en algún momento
choca, es imposible. Incluye la plantilla de recorrido sin recursión, necesaria en
Python para grafos grandes.
\ej un ciclo de 4 nodos se puede partir; uno de 3 no, porque tiene longitud impar.

\guia{Árbol de Expansión Mínima (Kruskal)}{alg:4-13}
Conectar todo al menor costo posible: cables, carreteras, tuberías. Se ordenan las
conexiones de más barata a más cara y se van tomando las que unen grupos todavía
separados. También sirve para agrupar: cortar las conexiones más caras del árbol
deja los grupos más separados entre sí.
\ej con tres ciudades y costos 1, 2 y 3, el mínimo para conectarlas es
\texttt{1 + 2 = 3}.

\guia{Puentes y árbol de puentes}{alg:4-14}
Qué conexiones son críticas: si esa calle se corta, la red se parte. Contrayendo
todo lo que sigue conectado sin ella, el grafo se convierte en un árbol, y sobre
ese árbol se responden preguntas que en el grafo original serían difíciles.
\ej para maximizar cuántos tramos obligatorios hay entre dos puntos, se toma el
camino más largo del árbol de puentes.

\guia{Orden Topológico + DP en DAG}{alg:4-15}
En un grafo sin ciclos, ordena las tareas de modo que nada venga antes que sus
requisitos. Recorriendo en ese orden, cualquier cálculo acumulado (camino más
largo, cuántos caminos hay, costo mínimo) sale en una sola pasada.
\ej la ruta que visita más ciudades de 1 a $n$, y la lista de ciudades por las que
pasa.

\guia{Bellman-Ford (ciclos negativos)}{alg:4-16}
Para grafos donde avanzar puede DAR beneficio en vez de costar, que es justo donde
Dijkstra se rompe. Además detecta el caso patológico: un ciclo que te deja mejor
cada vez que lo recorres, o sea, ganancia infinita.
\ej detectar arbitraje en cambio de divisas: una secuencia de cambios que te
devuelve más dinero del que pusiste.

\guia{LCA y distancias en árbol}{alg:4-17}
En un árbol genealógico o una jerarquía, el ancestro común más cercano de dos
nodos. Con él se obtiene al instante la distancia entre dos nodos cualesquiera, sin
recorrer el árbol en cada consulta.
\ej en un árbol de 200\,000 nodos, cada pregunta de distancia se responde de
inmediato tras un precálculo único.

\guia{Flujo Máximo y Corte Mínimo (Dinic)}{alg:4-18}
Cuánto se puede mandar de un punto a otro por una red con capacidades, y el
resultado gemelo: cuál es el conjunto más barato de conexiones a cortar para
separarlos del todo. Los dos números son el mismo, y ese es el teorema que hace
útil el algoritmo.
\ej si hay \texttt{2} rutas independientes entre el banco y el puerto, hacen falta
exactamente \texttt{2} cierres de calle para incomunicarlos.

\guia{2-SAT}{alg:4-19}
Cada elemento tiene dos opciones y hay restricciones de la forma "si esta, entonces
aquella". Decidir si existe una combinación que las cumpla todas es rapidísimo con
dos opciones, aunque con tres el mismo problema sea intratable. Se resuelve
reutilizando las componentes fuertemente conexas.
\ej con $n$ personas que piden cada una dos deseos sobre los ingredientes, dice si
hay una pizza que contente a todos.

% ============================================================
\subsection*{Búsqueda, Greedy y Optimización}

\guia{Búsqueda Binaria sobre la Respuesta}{alg:5-1}
El truco mental más rentable del repertorio: cuando calcular la respuesta es
difícil pero VERIFICAR un candidato es fácil, se adivina el candidato a la mitad y
se descarta media búsqueda cada vez. Funciona siempre que valga "si $X$ sirve,
todo lo mayor también".
\ej para la $k$-ésima suma más pequeña entre pares, se prueba un valor y se cuenta
cuántas sumas quedan por debajo.

\guia{Enumeración por Bitmask}{alg:5-2}
Cuando hay pocas decisiones de sí/no (unas 20 como mucho) y se pueden probar todas.
Cada número entero, leído en binario, representa una combinación distinta, así que
un simple bucle las recorre todas sin necesidad de recursión.
\ej con 16 bits desconocidos se prueban las \texttt{65\,536} combinaciones
posibles y se devuelve la primera válida.

\guia{Maximizar el mínimo}{alg:5-3}
La variante más común de la búsqueda binaria: repartir algo para que el peor
elemento quede lo mejor posible. "Maximizar el mínimo" y "minimizar el máximo" son
la misma idea, y casi siempre se resuelven así.
\ej con una muralla de alturas desiguales y un solo refuerzo disponible, la mayor
altura mínima que se puede garantizar.

\guia{Huffman (fusión greedy)}{alg:5-4}
Cuando hay que juntar (o partir) grupos repetidamente y cada operación cuesta el
tamaño del grupo resultante, lo óptimo es unir siempre los dos más pequeños. La
misma idea aparece disfrazada en "unir cuerdas", "mezclar archivos" y compresión de
datos.
\ej pesos \texttt{[1, 2, 3]}: unir 1+2 (cuesta 3) y luego 3+3 (cuesta 6) da un
total de \texttt{9}, que es el mínimo.

\guia{Barrido de Eventos (Line Sweep)}{alg:5-5}
Para preguntas sobre muchos intervalos a la vez: cuánta gente hay como máximo al
mismo tiempo, cuántas salas hacen falta, cuándo se solapan. Cada intervalo se
parte en un evento de entrada y uno de salida, se ordenan por tiempo y se recorren
llevando un contador. Nada de revisar par contra par.
\ej clientes que están de 5 a 8, de 2 a 4 y de 3 a 9: como máximo hay
\texttt{2} a la vez.

\guia{Ventana Deslizante (dos punteros)}{alg:5-6}
El tramo contiguo más largo (o más corto) que cumple una condición. Los dos bordes
de la ventana solo avanzan, así que todo se recorre una vez aunque parezca que hay
que probar cada par de extremos. Sirve para "sin repetidos", "suma al menos $x$" o
"a lo sumo $k$ valores distintos".
\ej en \texttt{[1, 2, 1, 3, 2, 7, 4, 2]} el tramo más largo sin repetir es
\texttt{1, 3, 2, 7, 4}, de largo \texttt{5}.

% ============================================================
\subsection*{Programación Dinámica}

\guia{DP de Particionamiento sobre Residuos}{alg:6-1}
Partir una lista en grupos minimizando un costo, cuando el costo de cada grupo
depende solo de su suma módulo un número pequeño. El estado deja de ser "dónde
corté" (miles de opciones) y pasa a ser "qué resto llevaba" (diez), que es lo que
vuelve viable el cálculo en Python.
\ej pagar una compra por partes cuando cada parte se redondea al múltiplo de 10 más
cercano: conviene elegir dónde cortar.

\guia{Distancia de Edición y LCS}{alg:6-2}
Cuánto se parecen dos textos: el mínimo de inserciones, borrados y sustituciones
para convertir uno en el otro. Es la base de los correctores ortográficos y de
\texttt{diff}. La versión bit-paralela mete una columna entera de la tabla en un
solo número entero, y por eso corre en Python donde la versión normal no llega.
\ej de \texttt{"casa"} a \texttt{"calle"} la distancia es \texttt{3}.

\guia{Subsecuencia Creciente Más Larga (LIS)}{alg:6-3}
Cuántos elementos se pueden tomar, en orden, de forma que cada uno sea mayor que el
anterior. La versión obvia compara todos los pares y es lenta; la buena guarda, para
cada longitud, el valor más chico con que puede terminar, y ubica cada número nuevo
con búsqueda binaria. Aparece disfrazada en apilar cajas o encajar sobres: se ordena
por una dimensión y se busca la LIS en la otra.
\ej en \texttt{[7, 3, 5, 3, 6, 2, 9, 8]} la respuesta es \texttt{4}
(\texttt{3, 5, 6, 9}).

\guia{Mochila 0/1}{alg:6-4}
Elegir objetos, cada uno una sola vez, para sacar el mayor valor sin pasarse de una
capacidad. El error típico es tomar primero lo más valioso: a veces conviene más
una combinación de cosas medianas. La versión de esta entrada empaqueta una fila
entera de la tabla en un solo entero, que es lo que la hace viable en Python.
\ej con presupuesto 10 y libros de precio 4, 8 y 5 que dan 5, 12 y 8 páginas, lo
mejor no es el de 12 páginas sino el de 4 más el de 5: \texttt{13} páginas.

\guia{Caminos en Grilla}{alg:6-5}
Cuántas formas hay de cruzar una cuadrícula moviéndose solo a la derecha o hacia
abajo, esquivando casillas prohibidas. Cada casilla suma lo que le llega de arriba
y de la izquierda. Es el DP más simple en dos dimensiones y la base de sus
variantes: costo mínimo, máximo de objetos recogidos, pasar por puntos obligados.
\ej en una grilla $3\times3$ vacía hay \texttt{6} caminos; si el centro es trampa,
quedan \texttt{2} (los que van por el borde).

\guia{DP en Árboles}{alg:6-6}
Resolver un problema sobre un árbol yendo de las hojas hacia la raíz: cada nodo
decide con lo que ya le informaron sus hijos, guardando solo unos pocos estados
(por ejemplo, "estoy emparejado" o "estoy libre"). Con esa misma plantilla salen el
emparejamiento máximo, el mayor grupo sin vecinos y la cobertura mínima. Se recorre
sin recursión, porque en Python un árbol profundo la desbordaría.
\ej en el camino $1-2-3-4$ se pueden elegir \texttt{2} aristas sin compartir
nodos; en una estrella de cuatro puntas, solo \texttt{1}.

\guia{DP de Dígitos}{alg:6-7}
Contar cuántos números de un rango enorme cumplen una condición sobre sus dígitos
(sin dígitos repetidos juntos, suma de dígitos múltiplo de algo, sin cierto
dígito). Es imposible recorrer hasta \ten{18} números, así que se construyen
dígito a dígito, recordando solo lo necesario: si ya se quedó por debajo del
límite y cuál fue el último dígito.
\ej entre \texttt{123} y \texttt{321} hay \texttt{171} números sin dos dígitos
adyacentes iguales.

\guia{DP de Intervalos}{alg:6-8}
Problemas donde se decide en qué orden partir o juntar un segmento, y el orden
cambia el costo. El estado es un tramo y se prueban todos sus puntos de corte. Así
se resuelven cortar barras, multiplicar cadenas de matrices y fusionar montones. Una
propiedad del costo (la optimización de Knuth) permite probar muchos menos puntos.
\ej en una barra de 100 con cortes en 25, 50 y 75, cortar primero en el centro
cuesta \texttt{200}; de izquierda a derecha, \texttt{225}.

\guia{DP con Ventana Mínima (deque monótona)}{alg:6-9}
Un DP de partición donde cada paso necesita el mínimo de los últimos $c$ valores.
Recalcularlo cada vez sería lento; una deque que conserva solo los candidatos útiles
lo da al instante, porque cada elemento entra y sale una sola vez. Es el uso más
frecuente de la ventana de tamaño fijo en concurso.
\ej con $c = 10$ y \texttt{[1, 1, 10, 10, 10, 10, 10, 9, 10, 10, 10, 10]}, lo
óptimo es separar \texttt{[1, 1]} del resto: costo total \texttt{92}.

% ============================================================
\subsection*{Matemáticas y Simulación}

\guia{Reducción Binaria (tipo Josefo)}{alg:7-1}
Procesos donde en cada ronda se elimina uno de cada dos y la cantidad se parte por
la mitad. Simular es imposible si hay miles de millones de elementos, pero como
cada ronda reduce el problema a la mitad, bastan unas pocas decenas de pasos.
\ej con mil millones de personas en fila, la respuesta sale en unos 30 pasos en
vez de mil millones.

\guia{Fuerza Bruta / Simulación Directa}{alg:7-2}
La entrada de reconocer cuándo NO hace falta ser listo: si los límites son
pequeños, hacer exactamente lo que dice el enunciado es más rápido de escribir y
mucho menos propenso a errores. Saber identificarlo ahorra minutos valiosos.
\ej cuántas personas hacen falta para que dos compartan cumpleaños: se van
multiplicando probabilidades hasta pasar el 50\%.

\guia{Simulación Lineal con Máximo de Prefijo}{alg:7-3}
Problemas que parecen complicados y se reducen a recorrer una vez y quedarse con el
máximo (o el mínimo). La dificultad está en darse cuenta de cuál es el cuello de
botella, no en programarlo.
\ej dos paredes que se acercan chocan en el tiempo que marca la fila más
estrecha.

\guia{Simulación de Eventos con Colas}{alg:7-4}
Sistemas donde llegan cosas en momentos distintos y hay un recurso compartido con
estado. En vez de avanzar segundo a segundo, se salta directamente al siguiente
evento relevante usando una cola por tipo.
\ej una escalera que solo va en un sentido a la vez: se calcula cuándo termina de
usarla la última persona.

\guia{Simulación XOR con Detección de Ciclos}{alg:7-5}
Procesos deterministas que pueden no terminar nunca. Si el estado se repite, es que
está dando vueltas y hay que cortar. Guardar los estados vistos es lo que
distingue "no hay solución" de un bucle infinito.
\ej interruptores que cambian grupos de bombillas: dice en cuántos pasos se apagan
todas, o que nunca ocurrirá.

\guia{Invariante de Paridad (ciclos de permutación)}{alg:7-6}
Para preguntas de "¿se puede llegar de esta configuración a esta otra?" donde
probar todas las combinaciones es imposible. Se busca una cantidad que NUNCA cambie
por más movimientos que hagas; si vale distinto en las dos configuraciones, la
respuesta es no y no hace falta buscar nada. En los rompecabezas deslizantes esa
cantidad es la paridad: cada movimiento intercambia dos piezas, y eso la invierte
siempre.
\ej el clásico 15-puzzle con las dos últimas piezas cambiadas es imposible de
resolver, y se sabe sin mover ni una ficha.

\guia{Fracciones Unitarias (Erdős--Straus)}{alg:7-7}
Escribir una fracción como suma de fracciones de numerador 1. El espacio de
búsqueda parece infinito, pero se puede acotar: la primera fracción no puede ser
demasiado pequeña, y eso limita cuánto hay que probar. Es el ejemplo típico de
"búsqueda que parece imposible hasta que demuestras una cota".
\ej $4/5 = 1/2 + 1/4 + 1/20$.

\guia{Fibonacci Rápido y Potencia de Matrices}{alg:7-8}
Calcular el término $n$ de una recurrencia cuando $n$ es astronómico. En vez de
avanzar de uno en uno, se duplica el índice en cada paso, así que llegar a \ten{18}
toma unos 60 pasos. La versión con matrices sirve para cualquier recurrencia lineal,
cambiando solo la matriz.
\ej $F_{10} = 55$; y $F_{10^{18}}$ módulo $\ten{9}+7$ sale en milisegundos.

% ============================================================
\subsection*{Construcción y Combinatoria}

\guia{Constructivo por Invariante}{alg:8-1}
Problemas de "construye algo que cumpla X" donde no hay que buscar nada: se propone
una fórmula y se argumenta que cierta propiedad (típicamente la paridad) nunca
puede romperse, sin importar las decisiones que se tomen.
\ej construir una permutación donde el resultado nunca pueda dar cero, sin
importar los signos que se elijan.

\guia{Constructivo Explícito}{alg:8-2}
El mismo tipo de problema resuelto dando directamente la respuesta a partir de una
identidad numérica conocida, en vez de programar una búsqueda.
\ej un arreglo cuya suma sea divisible por cada uno de sus elementos sale de la
identidad $1 = 1/2 + 1/3 + 1/6$.

\guia{Conteo por Simetría de Binomios}{alg:8-3}
Contar subconjuntos que cumplen una condición sin enumerarlos. La observación clave
es que, dentro de un grupo de elementos iguales, elegir una cantidad par o impar es
igual de probable, y eso colapsa el conteo a una potencia de dos.
\ej cuántas subsecuencias tienen suma alternada cero, con hasta \ten{5} elementos.

\guia{Greedy por Alternancia}{alg:8-4}
Construir una secuencia larga respetando límites de repetición consecutiva. Resulta
que alternar entre las dos opciones más holgadas siempre basta si el problema tiene
solución, así que no hace falta buscar.
\ej repartir noches entre varios alojamientos cuando cada uno tolera un máximo de
noches seguidas.

\guia{Stars and Bars}{alg:8-5}
Contar de cuántas formas se reparte una cantidad entre varios grupos, cuando los
elementos son indistinguibles y un grupo puede quedarse sin nada. Es una fórmula
cerrada, no un algoritmo.
\ej repartir \texttt{3} días de espera entre \texttt{2} paradas da \texttt{4}
maneras posibles.

% ============================================================
\subsection*{Cadenas}

\guia{Rotación Mínima (Booth)}{alg:9-1}
Cuando dos palabras deben considerarse iguales si una es la otra girada
(\texttt{"abc"} y \texttt{"cab"}). Da una forma canónica: la rotación
alfabéticamente menor. Con eso, comparar o eliminar duplicados es comparar dos
cadenas normales.
\ej \texttt{"cab"}, \texttt{"abc"} y \texttt{"bca"} colapsan todas a
\texttt{"abc"}.

\guia{Palíndromos por normalización}{alg:9-2}
El caso simple: comprobar si un texto se lee igual al derecho y al revés ignorando
espacios, mayúsculas y puntuación. Se limpia primero y se compara con su reverso;
no hace falta nada más sofisticado.
\ej \texttt{"Anita lava la tina"} es palíndromo una vez quitados espacios y
mayúsculas.

\guia{KMP sobre codificación}{alg:9-3}
Convierte un problema de árboles en uno de texto: si se escribe cada árbol como una
cadena en el orden correcto, "¿este subárbol tiene la misma forma que aquel?" se
vuelve "¿aparece esta cadena dentro de esta otra?", que se responde en una pasada.
\ej contar cuántos subárboles de una expresión tienen la misma forma que otra
expresión dada.

\guia{Suffix Array + LCP}{alg:9-4}
La herramienta pesada de cadenas: ordena alfabéticamente todos los finales posibles
del texto, y guarda cuánto comparte cada uno con el siguiente. Con eso salen la
subcadena repetida más larga, cuántas subcadenas distintas hay y la parte común
entre dos textos, sin revisar todas las subcadenas una por una.
\ej en \texttt{"banana"}, la subcadena repetida más larga es \texttt{"ana"}, de
longitud \texttt{3}.

\guia{Firma Canónica + Hash}{alg:9-5}
Agrupar cosas "equivalentes" sin compararlas todas contra todas. Se define una
huella que sale igual para todos los equivalentes (por ejemplo, las letras
ordenadas) y se usa como llave de un diccionario: cada consulta es un acceso
directo.
\ej \texttt{"roma"}, \texttt{"amor"} y \texttt{"mora"} comparten la huella
\texttt{"amor"}, así que caen en el mismo grupo.

\guia{Expresiones Regulares}{alg:9-6}
Cuando el problema se reduce a contar o localizar patrones fijos, el módulo
\texttt{re} ya está escrito y optimizado. La única trampa: por defecto no cuenta
apariciones que se solapan, y hay que usar una construcción especial si el patrón
puede solaparse consigo mismo.
\ej en \texttt{"aaaa"}, el conteo normal de \texttt{"aa"} da \texttt{2}, pero las
apariciones reales son \texttt{3}.

\guia{Subcadenas distintas con restricción}{alg:9-7}
Contar cuántos fragmentos DISTINTOS del texto cumplen una condición, sin contar dos
veces los repetidos. Combina el suffix array con una ventana que avanza, y es la
generalización natural de "cuántas subcadenas distintas tiene este texto".
\ej cuántos fragmentos distintos contienen a lo sumo $k$ letras "malas".

\guia{Z-Function}{alg:9-8}
Para cada posición dice cuánto coincide el texto consigo mismo empezando ahí. Es
prima de KMP y resuelve lo mismo, pero responde una pregunta más directa: "¿cuánto
del principio del texto vuelve a aparecer aquí?". Por eso es más cómoda en todo lo
que trata de prefijos —bordes, repeticiones, comparar el texto con partes de sí
mismo— mientras que KMP encaja mejor cuando buscas un patrón fijo dentro de un
texto largo. Muchos la prefieren simplemente porque es más corta de escribir y más
difícil de equivocar.
\ej en \texttt{"aabaab"} da \texttt{[-, 1, 0, 3, 1, 0]}: el \texttt{3} de la
cuarta posición avisa que \texttt{"aab"} vuelve a empezar ahí. Con entrada
\texttt{"fixprefixsuffix"} la salida es \texttt{"fix"}: es principio, es final, y
además aparece en medio.

\guia{Período Mínimo (función de prefijos)}{alg:9-9}
Detecta si una cadena es un bloque repetido y cuántas veces se repite. Sale de una
sola resta sobre la tabla de KMP, y es de los trucos con mejor relación entre lo
poco que cuesta y lo seguido que aparece.
\ej entrada \texttt{"abcabcabc"} $\rightarrow$ salida \texttt{3}; entrada
\texttt{"abcd"} $\rightarrow$ salida \texttt{1}.

\guia{Manacher}{alg:9-10}
El palíndromo más largo dentro de un texto, en una sola pasada. Probar todos los
centros y expandir sería demasiado lento con textos de un millón de caracteres;
Manacher aprovecha lo ya calculado para no repetir comparaciones. También sirve
para contar todos los fragmentos palindrómicos.
\ej entrada \texttt{"aybabtu"} $\rightarrow$ salida \texttt{"bab"}.

\guia{Aho-Corasick}{alg:9-11}
Buscar MUCHOS patrones a la vez en un texto, recorriéndolo una sola vez. Correr KMP
por separado para cada patrón sería lentísimo si hay miles; aquí todos los patrones
se meten en una misma estructura y el texto se lee de corrido. Es lo que hay detrás
de un filtro de palabras o un antivirus por firmas.
\ej con el texto \texttt{"abcab"} y los patrones \texttt{"ab"}, \texttt{"bc"},
\texttt{"xyz"}, la salida es \texttt{2}, \texttt{1}, \texttt{0}.
